\chapter{Testing Framework}

Testing the code you write is probably just as important as actually writing the code.
This is especially true when you are working on a complex system with many interacting components like an Operating System.

We therefore want to start by showing the Testing Framework we created for this project.
It proved useful throughout the different chapters of the course because it greatly simplified and standardized the way we would write tests.

% We created some macro magic that lets you create and run tests by doing this:
It uses some "macro magic" to create and group different test cases together.
These groups of commands can then be executed by using the \mintinline{c}{RUN_TESTS(<group>)} command.
It shows you if the tests succeeded or failed, and, if they failed, where and why they failed.

\begin{minted}{c}
CREATE_TEST(test_name1, group_name, 
{
    errval_t err;
    err = ...
    TEST_REQUIRE_OK(err);
    err = ...;
    TEST_REQUIRE_FAIL_WITH(err, SOME_ERR);
})

CREATE_TEST(test_name2, group_name, 
{
    // This makes a test fail but doesn't abort the execution.
    TEST_CHECK(1 == 2);
    // This makes a test fail and aborts execution of this test (but not the group).
    TEST_REQUIRE(1 == 2); 
})

// ...

void run_testing(void) {
    RUN_TESTS(group_name);
}
\end{minted}
This is how it looks like when a test fails (in red):
\begin{minted}{sh}
/lib/grading/group_name_tests.h:12: 1==2 failed
/lib/grading/group_name_tests.h:13: 1==2 failed
Test test_name2 in group_name failed.
1/8 tests in group_name failed.
\end{minted}
This is how it looks like when all tests pass (in green):
\begin{minted}{sh}
All 8/8 tests in group_name succeeded.
\end{minted}